%% kapa: exam=1 round=34 year=2023 period=2 form=A subject=appraisal_law subject_ko=감정평가관계법규 seq=29 number=29
%% source: Q-Net 감정평가사 기출문제 (cst00309) artlSeq=5211789 file=raw/1008/5211789/2023년도 제34회 감정평가사 1차시험 2교시.pdf
%% conversion: kordoc 4.24.51 -> tools/gen1.py (auto); review=proofread-sonnet; flags=none
\begin{question}{29}
공간정보의 구축 및 관리 등에 관한 법령상 지적에 관한 설명으로 옳지 않은 것은?
\begin{choices}
  \item 지번은 지적소관청이 지번부여지역별로 차례대로 부여한다.
  \item 면적의 단위는 제곱미터로 한다.
  \item 지적도면의 번호는 경계점좌표등록부에 등록하여야 할 사항에 속한다.
  \item 토지소유자는 신규등록할 토지가 있으면 그 사유가 발생한 날부터 90일 이내에 지적소관청에 신규등록을 신청하여야 한다.
  \item 행정구역의 명칭이 변경되었으면 지적공부에 등록된 토지의 소재는 새로운 행정구역의 명칭으로 변경된 것으로 본다.
\end{choices}
\end{question}
\answer{4}
